\documentclass[../main.tex]{subfiles}
\begin{document}

\chapter{Multi-Processing}
\lessoninfo{
  video={Lesson 18, videos 1--22},
  textbook={H\&P \cite{hennessy2017} §1.1--1.2, §1.5, §5.1, ch.~3, Appendix B},
  keywords={Flynn's taxonomy, MIMD, multicore processor, centralized shared memory, UMA, SMP, distributed shared memory, NUMA, message passing, shared memory, chip multiprocessor, fine-grained multithreading, simultaneous multithreading, working set},
}

Up to this point the course has made a single processor run a single
program as fast as possible---with pipelining, branch prediction,
out-of-order execution and a memory hierarchy---and Lesson 17 made such a
system dependable.  This lesson turns to machines that execute several
programs, or several parts of one program, at the same time.  It classifies
parallel machines, explains why the additional transistors of each
technology generation now become additional cores rather than a faster
single core, compares the two ways of connecting processors to memory and
the two programming models that go with them, and ends with multithreading
hardware, in particular simultaneous multithreading, which lets one core
execute several threads at once.

\section{Flynn's taxonomy}
\label{sec:l18-flynn}

A machine can do more than one thing at a time in different
ways.\source{Lesson 18, videos 1--2; H\&P §1.2; \cite{flynn1966}.}  The
classic classification of parallel machines,
\term{Flynn's taxonomy}[Flynn の分類], looks at two numbers: how many \emph{instruction streams}---each a
sequence of instructions with its own program counter---the machine executes
at the same time, and how many \emph{data streams} these instructions
operate on.  Distinguishing \enquote{one} from \enquote{more than one} for
each number gives four classes (\cref{tab:l18-flynn}).

\begin{table}[htbp]
  \centering\small
  \caption{Flynn's taxonomy of parallel machines.}
  \label{tab:l18-flynn}
  \begin{tabular}{@{}lccl@{}}
    \toprule
    Class & Instruction streams & Data streams & Example \\
    \midrule
    SISD & 1        & 1        & uniprocessor \\
    SIMD & 1        & $>1$     & vector processor, SSE/AVX \\
    MISD & $>1$     & 1        & stream processor (rare) \\
    MIMD & $>1$     & $>1$     & multiprocessor, multicore \\
    \bottomrule
  \end{tabular}
\end{table}

\begin{description}
  \item[SISD] A \term{single instruction, single data}[SISD]%
    \index{SISD|see{single instruction, single data}} machine executes one
    instruction stream on one stream of data: the uniprocessor of the
    previous lessons.  A superscalar out-of-order core is still SISD; the
    instructions it executes in the same cycle all belong to one program.
  \item[SIMD] A \term{single instruction, multiple data}[SIMD]%
    \index{SIMD|see{single instruction, multiple data}} machine executes one
    instruction stream, but each instruction operates on several data items.
    Vector processors are of this kind, and so are the multimedia extensions
    of current instruction sets (MMX, SSE and AVX in x86), whose
    instructions apply one operation to every element of a short vector.
  \item[MISD] A \term{multiple instruction, single data}[MISD]%
    \index{MISD|see{multiple instruction, single data}} machine executes
    several instruction streams on one data stream; a stream processor in
    which several programs each process the same incoming data would be an
    example.  Such machines are rare.
  \item[MIMD] A \term{multiple instruction, multiple data}[MIMD]%
    \index{MIMD|see{multiple instruction, multiple data}} machine executes
    several instruction streams, each on its own data.
\end{description}

A machine that consists of several processors, each executing its own
program (or its own part of a program) on its own data, is a
\term{multiprocessor}[マルチプロセッサ]; multiprocessors are MIMD machines,
and almost every parallel machine in use today belongs to this class.

\section{Why multiprocessors}
\label{sec:l18-why}

A uniprocessor can be made faster in two ways: by completing more
instructions per cycle or by making the cycles shorter.\source{Lesson 18,
videos 3--4; H\&P §1.1, §1.5.}  Both have run out of room.
\begin{itemize}
  \item \textbf{Wider.}  Current cores already handle about four
    instructions per cycle.  Programs seldom contain enough independent
    instructions to fill even more slots, so additional width yields
    diminishing returns (Lesson 6).
  \item \textbf{Faster.}  A higher clock frequency requires a higher supply
    voltage.\margin{Clock rates stopped rising instead: the Pentium~4
    reached \SI{3.8}{\giga\hertz} in 2004, and the fastest desktop parts
    today reach about \SI{6}{\giga\hertz}.  H\&P §1.5.}  Since dynamic
    power is proportional to $V^{2}f$
    (\cref{eq:active-power}), raising $f$ and $V$ together makes power grow
    roughly with $f^{3}$, and processors already dissipate as much power as
    can be removed at reasonable cost.
\end{itemize}

Moore's law, on the other hand, continues: the number of transistors that
fit on a chip still doubles about every 18 months.  Instead of a more
complex single core, the transistors can be used to replicate the core.  A
\term{multicore processor}[マルチコアプロセッサ] is a single chip that
contains two or more complete processors, called \emph{cores}.  Doubling
the number of cores in every generation can double the performance in every
generation---provided that the software uses all the cores.\caution{Moore's
own figure was one year in 1965 \cite{moore1965} and two years in his 1975
revision; \enquote{18 months} is a later popularization.}

\begin{example}
  \label{ex:l18-power}
  A core runs at \SI{3}{\giga\hertz} and dissipates \SI{40}{\watt} of dynamic
  power.  Running it at \SI{3.6}{\giga\hertz}, $1.2\times$ faster, requires
  a proportionally higher voltage and raises the power by a factor of about
  $1.2^{3}\approx 1.73$, to \SI{69}{\watt}.  Two copies of the original core
  at \SI{3}{\giga\hertz} dissipate \SI{80}{\watt} and complete twice the
  work of one core---if the program keeps both busy.  Two copies at
  \SI{2.4}{\giga\hertz} ($0.8\times$) dissipate
  $2\times\SI{40}{\watt}\times 0.8^{3}\approx\SI{41}{\watt}$, about as much
  as the single original core, yet together complete $1.6\times$ its work.
\end{example}

\section{Multiprocessors need parallel programs}
\label{sec:l18-parallel-programs}

A faster uniprocessor speeds up every program, including programs written
long ago.  A multiprocessor speeds up only programs that divide their work
among several cores.\source{Lesson 18, video 5; H\&P §5.1.}  Compared with
sequential (single-threaded) code, such parallel programs are harder in
every respect:
\begin{itemize}
  \item \textbf{Development.}  Besides solving the problem, the programmer
    must divide the work among threads and coordinate them.
  \item \textbf{Debugging.}  The behavior of a parallel program depends on
    the relative timing of its threads, which changes from run to run, so an
    error may appear in one run and not in the next.
  \item \textbf{Performance scaling.}  Ideally, performance grows linearly
    with the number of cores.  In practice it grows up to some number of
    cores and then levels off; improving the program raises the level, but
    performance still stops growing at some point (\cref{fig:l18-scaling}).
\end{itemize}

\begin{marginfigure}
  \resizebox{\linewidth}{!}{% Speedup versus number of cores: ideal linear scaling, and the levelling
% off predicted by Amdahl's law for parallel fractions 0.97 and 0.9.
\begin{tikzpicture}[x=0.118cm,y=0.118cm, font=\scriptsize\sffamily]
  \draw[ln-rule] (0,0) grid[xstep=8,ystep=8] (32,32);
  \draw[->] (0,0) -- (34,0);
  \draw[->] (0,0) -- (0,34);
  \foreach \v in {0,8,16,24,32} {
    \node[below, inner sep=1.5pt] at (\v,0) {\v};
    \node[left, inner sep=1.5pt] at (0,\v) {\v};
  }
  \node[below] at (16,-3.2) {\iflangja{コア数 $N$}{cores $N$}};
  \node[rotate=90, above] at (-4.6,16) {\iflangja{高速化率}{speedup}};
  \draw[thick, ln-muted, dashed] (1,1) -- (32,32)
    node[pos=0.78, above left, inner sep=1pt] {\iflangja{理想}{ideal}};
  \draw[thick, ln-accent] plot[domain=1:32, samples=48, smooth]
    (\x, {1/(0.03+0.97/\x)});
  \node[ln-accent, above left, inner sep=1pt] at (32,16.6) {$f_{\mathrm{par}}=0.97$};
  \draw[thick, ln-caution] plot[domain=1:32, samples=48, smooth]
    (\x, {1/(0.1+0.9/\x)});
  \node[ln-caution, above left, inner sep=1pt] at (32,7.8) {$f_{\mathrm{par}}=0.9$};
\end{tikzpicture}
}
  \caption{Speedup on $N$ cores: ideal scaling and
    \cref{eq:l18-amdahl-cores}.}
  \label{fig:l18-scaling}
\end{marginfigure}
Amdahl's law (\cref{eq:amdahl}) shows why the levelling off is
unavoidable.  If a fraction $f_{\mathrm{par}}$ of a program's execution
time is divided evenly among $N$ cores and the rest remains sequential, the
speedup is at most
\begin{equation}
  S(N) \;=\; \frac{1}{(1-f_{\mathrm{par}}) + \dfrac{f_{\mathrm{par}}}{N}}
  \;<\; \frac{1}{1-f_{\mathrm{par}}} .
  \label{eq:l18-amdahl-cores}
\end{equation}
Coordinating the threads costs additional time on top of this, so real
programs stay below these curves.\tip{Half the ideal speedup, $S(N)=N/2$,
needs $f_{\mathrm{par}}=(N-2)/(N-1)$: \SI{86}{\percent} on 8 cores,
\SI{98}{\percent} on 64.}

\begin{example}
  In a program, \SI{95}{\percent} of the time is spent in a loop whose
  iterations are independent, so $f_{\mathrm{par}}=0.95$.  On 8 cores the speedup is at
  most $1/(0.05+0.95/8)\approx 5.9$, on 32 cores
  $1/(0.05+0.95/32)\approx 12.5$, and no number of cores takes it beyond
  $1/0.05 = 20$.
\end{example}

\needspace{5\baselineskip}
\begin{keypoint}
  Replicating cores turns transistors into throughput far more efficiently
  than building a larger or faster single core, but only for programs that
  are written to use many cores, and only to the extent that their
  sequential part allows.
\end{keypoint}

\section{Centralized shared memory}
\label{sec:l18-centralized}

The simplest multiprocessor connects all processors to one main
memory.\source{Lesson 18, videos 6--8; H\&P §5.1.}

\begin{definition}[Centralized shared memory]
  In a multiprocessor with \term{centralized shared memory}[集中共有メモリ]
  all cores, each with its own caches, are connected by one bus to a single
  main memory and to the I/O devices (\cref{fig:l18-centralized}).
\end{definition}

\begin{figure}[htbp]
  \centering
  % Centralized shared memory (UMA, SMP): cores with private caches share one
% main memory and the I/O devices through a common bus.
\begin{tikzpicture}[x=1cm,y=1cm, font=\footnotesize\sffamily,
  core/.style={draw, fill=ln-accent!15, minimum width=1.7cm, minimum height=0.6cm, inner sep=2pt},
  cache/.style={draw, fill=ln-box, minimum width=1.7cm, minimum height=0.5cm, inner sep=2pt},
  big/.style={draw, fill=ln-box, minimum height=0.7cm, inner sep=2pt}]
  \foreach \i in {0,...,3} {
    \node[core] (c\i) at (\i*2.3,0) {\iflangja{コア}{core}~\i};
    \node[cache] (k\i) at (\i*2.3,-0.85) {\iflangja{キャッシュ}{cache}};
    \draw (c\i.south) -- (k\i.north);
    \draw (k\i.south) -- (\i*2.3,-1.7);
  }
  \draw[line width=1.4pt] (-1.0,-1.7) -- (7.9,-1.7)
    node[right, font=\footnotesize] {\iflangja{バス}{bus}};
  \node[big, minimum width=4.4cm] (mem) at (1.15,-2.75) {\iflangja{主記憶}{main memory}};
  \node[big, minimum width=2.8cm] (io) at (5.75,-2.75) {\iflangja{I/O 装置}{I/O devices}};
  \draw (mem.north) -- (1.15,-1.7);
  \draw (io.north) -- (5.75,-1.7);
\end{tikzpicture}

  \caption{Centralized shared memory: all cores reach the main memory and
    the I/O devices through the same bus.}
  \label{fig:l18-centralized}
\end{figure}

Any core can read and write any memory location and use any I/O device, so
cores exchange data simply by writing it to memory where another core reads
it.  Because the memory is equally far from all cores, every core needs the
same time to access a given location: the machine has
\term{uniform memory access}[UMA]\index{UMA|see{uniform memory access}}
(UMA).  And because no core has a special role, the machine is also called
a \term{symmetric multiprocessor}[対称型マルチプロセッサ]%
\index{SMP|see{symmetric multiprocessor}} (SMP).  A multicore chip whose
cores share the main memory attached to it is the most common example.

Centralized shared memory does not scale to many cores, for two reasons.
\begin{itemize}
  \item \textbf{Memory size.}  One memory holds the data of all cores, so it
    must grow with the number of cores, and a larger memory is a slower
    memory (Lessons 12 and 15).
  \item \textbf{Memory bandwidth.}  The cache misses of all cores and all
    I/O transfers go over the same bus.  The more cores there are, the more
    they wait for each other's memory accesses, until memory access is
    effectively serialized and adding cores brings no further benefit.
\end{itemize}
Centralized shared memory therefore works well only for small machines,
with about 2 to 16 cores.

\section{Distributed shared memory}
\label{sec:l18-distributed}

Larger machines distribute the memory among the
processors.\source{Lesson 18, videos 9--10; H\&P §5.1.}

\begin{definition}[Distributed shared memory]
  In a multiprocessor with \term{distributed shared memory}[分散共有メモリ]
  each core has its own local memory and a network interface, and the cores
  are connected by a network (\cref{fig:l18-distributed}).  A core accesses
  only its local memory directly; data held by another core is obtained by
  sending that core a request over the network.
\end{definition}

\begin{figure}[htbp]
  \centering
  % Distributed shared memory (NUMA): every node has its own memory and a
% network interface; memory of other nodes is reached through the network.
\begin{tikzpicture}[x=1cm,y=1cm, font=\footnotesize\sffamily, >={Stealth[length=1.8mm]},
  core/.style={draw, fill=ln-accent!15, minimum width=2.1cm, minimum height=0.55cm, inner sep=1pt},
  part/.style={draw, fill=ln-box, minimum width=2.1cm, minimum height=0.5cm, inner sep=1pt},
  small/.style={draw, fill=ln-box, minimum width=0.95cm, minimum height=0.5cm, inner sep=1pt, font=\scriptsize\sffamily}]
  \foreach \i in {0,...,3} {
    \begin{scope}[xshift=\i*2.75cm]
      \draw[ln-rule, rounded corners=2pt] (-1.25,0.5) rectangle (1.25,-2.75);
      \node[core] (c\i) at (0,0) {\iflangja{コア}{core}~\i};
      \node[part] (k\i) at (0,-0.8) {\iflangja{キャッシュ}{cache}};
      \node[small] (m\i) at (-0.575,-2.2) {\iflangja{メモリ}{memory}};
      \node[small] (n\i) at (0.575,-2.2) {NIC};
      \draw (c\i.south) -- (k\i.north);
      \draw (k\i.south) -- (0,-1.45);
      \draw (-0.575,-1.45) -- (0.575,-1.45);
      \draw (m\i.north) -- (-0.575,-1.45);
      \draw (n\i.north) -- (0.575,-1.45);
      \draw (n\i.south) -- (0.575,-3.3);
    \end{scope}
  }
  \draw[line width=1.4pt] (-1.25,-3.3) -- (9.5,-3.3);
  \node[font=\footnotesize, anchor=north west] at (-1.25,-3.35) {\iflangja{ネットワーク}{network}};
  % local access of core 0 to its own memory
  \draw[->, thick, ln-tip] (-0.12,-1.05) -- (-0.12,-1.33) -- (-0.42,-1.33) -- (-0.42,-1.95);
  \node[font=\scriptsize, text=ln-tip, anchor=east, inner sep=1pt] at (-0.64,-1.7)
    {\iflangja{高速}{fast}};
  % remote access from core 0 to the memory of node 2, through both NICs
  \begin{scope}[dashed, thick, ln-caution]
    \draw (0.12,-1.05) -- (0.12,-1.33) -- (0.8,-1.33) -- (0.8,-1.95);
    \draw (0.8,-2.45) -- (0.8,-3.05) -- (6.3,-3.05) -- (6.3,-2.45);
    \draw[->] (6.3,-1.95) -- (6.3,-1.33) -- (5.2,-1.33) -- (5.2,-1.95);
  \end{scope}
  \node[font=\scriptsize, text=ln-caution, anchor=north] at (3.4,-3.35)
    {\iflangja{他ノードのメモリ：ネットワーク経由（低速）}{memory of another node: via the network (slow)}};
\end{tikzpicture}

  \caption{Distributed shared memory: each node combines a core, its cache, local
    memory and a network interface card (NIC).}
  \label{fig:l18-distributed}
\end{figure}

A core with its memory and network interface is almost a computer of its
own, a \emph{node}; such a machine is therefore also called a
\emph{cluster} or a \emph{multicomputer}.  A core reaches its local memory
as fast as a uniprocessor does, but data held by another node must travel
over the network, which takes far longer.  The access time depends on where
the data is, so the machine has
\term{non-uniform memory access}[NUMA]%
\index{NUMA|see{non-uniform memory access}} (NUMA).

Distributed shared memory scales to very large numbers of cores; clusters and
supercomputers are built this way.  The reason is not that the network is
fast---it is much slower than a memory bus---but that no memory or bus is
shared by all cores, and that the programmer, who has to request every
remote access explicitly, is forced to think about communication and
therefore keeps it to a minimum.\margin{Roughly: \SI{80}{\nano\second} to
local DRAM, \SI{130}{\nano\second} to the memory of the other socket,
1--\SI{2}{\micro\second} to another node over InfiniBand.}

Since a remote access costs far more than a local one, placement of data
decides the performance of a NUMA machine.  The operating system should
place each page in the memory of the node whose core accesses it most---the
private data of a thread, such as its stack, on the node where the thread
runs---and keep the thread on that node, so that most memory accesses stay
local.  A page that several cores use, however, cannot be local to all of
them.\tip{On a NUMA machine, keep each thread on the node that holds the
data it uses most.  Linux places a page on the node whose thread first
touches it; \texttt{numactl} overrides the choice.}

\begin{note}[Terminology]
  Following the lecture, \enquote{distributed shared memory} is used here
  for every multiprocessor whose memory is divided among its nodes, although
  with explicit requests over a network no memory is actually shared.  H\&P
  reserve the name for machines whose nodes do share one physical address
  space---any core can issue any address, and the interconnect carries an
  access to a remote address to the node that holds it---and call machines
  that communicate only by explicit messages clusters.  The lecture calls
  both kinds NUMA; in H\&P's terminology NUMA refers only to the machines
  with a shared address space, in which the access time depends on where the
  data is.
\end{note}

\section{Message passing and shared memory}
\label{sec:l18-programs}

A parallel program can be written in two fundamentally different
ways.\source{Lesson 18, videos 11--14; H\&P §5.1.}

\begin{definition}[Message passing]
  In the \term{message passing}[メッセージパッシング] programming model
  each process has its own memory, which no other process can access.
  Processes exchange data only by explicitly sending and receiving messages.
\end{definition}

\begin{definition}[Shared memory]
  In the \term{shared memory}[共有メモリ] programming model all threads
  share one address space and exchange data simply by writing variables
  that other threads read.
\end{definition}

The two models are compared here on the same task: adding up the $N$
elements of an array with $P$ processes or threads.  In both programs every
process or thread executes the same code and differs only in its
identifier \texttt{myPID}, which ranges from 0 to $P-1$.

\subsection{A message-passing program}

Each process holds only the $N/P$ elements assigned to it.  Distributing
the parts of the array is itself explicit work (not shown): the process
that initially has the data must send each other process its part.  Each
process then adds up its own part, and process 0 receives and combines the
partial sums of all the others (\cref{ex:l18-mp}).

\needspace{8\baselineskip}
\begin{example}[Message passing]
  \label{ex:l18-mp}
  Every process executes the following code, with its own $N/P$ elements
  already in \texttt{part}:
  \begin{lstlisting}[language=C, numbers=none]
#define N 4096            /* array size          */
#define P 8               /* number of processes */
double part[N/P];         /* this process's part */
double mySum = 0;
for (int i = 0; i < N/P; i++)
  mySum += part[i];
if (myPID == 0) {
  for (int p = 1; p < P; p++) {
    double pSum;
    recv(p, &pSum);       /* wait for p's sum    */
    mySum += pSum;
  }
  printf("sum: %f\n", mySum);
} else
  send(0, &mySum);        /* to process 0        */
\end{lstlisting}
\end{example}

\subsection{A shared-memory program}

With shared memory the whole array is visible to every thread,\margin{In
practice these are libraries: MPI for message passing---%
\texttt{MPI\_Reduce} is exactly this sum---and OpenMP or pthreads for
shared memory.} so nothing
has to be distributed: each thread computes the range of indices it is
responsible for and adds its partial sum to a shared total
(\cref{ex:l18-sm}).  Variables declared \texttt{shared} exist once for all
threads; all other variables are private to each thread.

\begin{example}[Shared memory]
  \label{ex:l18-sm}
  Every thread executes the following code on the shared array:\caution{%
  \texttt{shared} is notation, not C\@.  With pthreads, globals and the
  heap are shared and the stack is private.}
  \begin{lstlisting}[language=C, numbers=none, morekeywords={shared}]
#define N 4096
#define P 8
shared double array[N];
shared double allSum = 0;
shared lock_t sumLock;
double mySum = 0;         /* private to the thread */
int lo = myPID * N / P;
int hi = (myPID + 1) * N / P;
for (int i = lo; i < hi; i++)
  mySum += array[i];
lock(&sumLock);
allSum += mySum;
unlock(&sumLock);
if (myPID == 0)           /* see text              */
  printf("sum: %f\n", allSum);
\end{lstlisting}
\end{example}

The lock is essential.  The statement \texttt{allSum += mySum} reads
\texttt{allSum}, adds, and writes the result back; if two threads both read
the old value before either writes, one of the two partial sums is lost.
The lock admits one thread at a time to the update.\tip{Combining $P$
partial sums one after another takes $P-1$ steps; in a tree it takes
$\log_{2}P$.}  The program is also
still incomplete: thread 0 must wait until all other threads have added
their sums before it prints the total.  Locks and this kind of waiting are
the subject of Lesson 20.

\subsection{Comparison}

\Cref{tab:l18-mp-sm} summarizes how the two models differ.  In message
passing, every transfer of data and the decision which process holds which
data are written by the programmer, and the hardware only has to connect
otherwise independent nodes by a network.  With shared memory,
communication and data distribution happen implicitly through the shared
variables, but the hardware must provide a shared address space, including
consistent contents in the caches of all cores (Lesson 19), and support for
synchronization (Lesson 20).  The two models also require explicit code in
different places: message passing needs code to distribute the data, but a
receive already waits until the data has arrived, whereas shared memory
needs no distribution code but explicit synchronization wherever a thread
must wait for the others.\caution{That a receive waits is also a trap: if
two processes each send to the other before receiving, both can block for
ever.}

\begin{table}[htbp]
  \centering\small
  \caption{Message passing and shared memory compared.}
  \label{tab:l18-mp-sm}
  \begin{tabular}{@{}lll@{}}
    \toprule
                                   & Message passing   & Shared memory \\
    \midrule
    Communication                  & by the programmer & automatic \\
    Data distribution              & manual            & automatic \\
    Hardware support               & simple            & extensive \\
    Programming for correctness    & difficult         & less difficult \\
    Programming for performance    & difficult         & very difficult \\
    \bottomrule
  \end{tabular}
\end{table}

For the programmer the two models distribute the difficulty differently.
A shared-memory program is easier to get right: a sequential program can be
parallelized one loop at a time, with locks around the updates of shared
data.  Making it fast is much harder, because the costs of sharing do not
appear in the source code---threads waiting for locks, data moving between
the caches of different cores, remote accesses on a NUMA machine.  A
message-passing program is harder to get right, since the data distribution
and every send with its matching receive must be written by hand, but once
it is correct all communication is visible in the code, and much of the
work of making it fast has been done along the way.

\begin{keypoint}
  Message passing makes communication explicit: hard to program correctly,
  but its costs are visible.  Shared memory makes communication implicit:
  easier to program correctly, but much harder to make fast.
\end{keypoint}

\section{Shared-memory hardware}
\label{sec:l18-mt-hardware}

The threads of a shared-memory program can be run by the hardware in
several ways.\source{Lesson 18, video 15; H\&P ch.~3, §5.1.}
\begin{enumerate}
  \item \textbf{Several cores sharing a physical address space.}  Every core
    can issue any physical address.  The memory may be centralized (UMA) or
    distributed among the nodes, which then share one physical address space
    (NUMA; see the note in \cref{sec:l18-distributed}).  Each thread runs on a
    core of its own, and the threads truly execute at the same time.
  \item \textbf{One core, time-sharing.}  The operating system runs one
    thread for a while, then saves its state (PC and registers) and restores
    the state of another thread.  Such \emph{context switches} need no
    special hardware, and the threads automatically use the same physical
    memory, but only one thread executes at any moment: time-sharing gains
    nothing in performance.
  \item \textbf{One core, hardware multithreading.}  The core itself keeps
    the state of several threads and switches among them without the
    operating system.  By how often the core changes threads, three
    variants are distinguished, each requiring more hardware support than
    the one before:
    \begin{itemize}
      \item \term{coarse-grained multithreading}[粗粒度マルチスレッディング]
        changes to another thread only after several cycles, typically when
        the running thread has to wait for a slow event such as a miss in
        the last-level cache;
      \item \term{fine-grained multithreading}[細粒度マルチスレッディング]
        can change to another thread in every cycle;
      \item \term{simultaneous multithreading}[同時マルチスレッディング]%
        \index{SMT|see{simultaneous multithreading}} (SMT) executes
        instructions from several threads in the same cycle.  Intel's
        Hyper-Threading is an implementation of SMT.
    \end{itemize}
\end{enumerate}

\section{Multithreading performance}
\label{sec:l18-mt-performance}

Consider two threads and a superscalar core that can handle four
instructions per cycle.\source{Lesson 18, videos 16--17; H\&P ch.~3;
\cite{olukotun1996,tullsen1995}.}  A single thread rarely fills all four
slots of a cycle: dependences among its instructions, mispredicted branches
and cache misses leave many slots empty, and in some cycles the thread
cannot proceed at all.  \Cref{fig:l18-mt-timing} shows schematically how
four configurations use the slots.  In all four, an instruction for which
no slot is free waits for the next cycle, and a thread that waits for an
event, such as a cache miss, waits the same number of cycles whether or
not it holds the core.

\begin{figure}[htbp]
  \centering
  % Schematic use of the four slots per cycle of a core by two threads over
% ten cycles: (a) time-sharing, (b) chip multiprocessor, (c) fine-grained
% multithreading, (d) simultaneous multithreading.
% Alone, thread 1 executes 3,2,3,-,-,1,1,-,4,2 and thread 2
% 3,1,2,-,2,-,-,4,1,3 instructions in the ten cycles (16 each; "-" is a
% cycle spent waiting for an event such as a cache miss).  One model for all
% panels: instructions for which no slot is free wait for the next cycle
% (with SMT, 4 slots are split 2+2 when both threads have more ready), and
% a wait elapses whether or not the thread holds the core.  (a) switches
% threads after five cycles; (c) alternates between the threads that are
% not waiting.
\begin{tikzpicture}[x=1cm,y=1cm, font=\scriptsize\sffamily]
  \def\lncw{0.24}
  % \lnslots{x0}{y0}{cycle/slots of thread 1/slots of thread 2, ...}
  \def\lnslots#1#2#3{%
    \foreach \c/\a/\b in {#3} {
      \foreach \s in {1,...,4} {
        \pgfmathtruncatemacro\lnk{\s<=\a ? 1 : (\s<=\a+\b ? 2 : 0)}
        \ifcase\lnk \def\lnf{white}\or \def\lnf{ln-accent!60}\or \def\lnf{ln-tip!45}\fi
        \draw[ln-rule, fill=\lnf] ({#1+(\c-1)*\lncw},{#2-(\s-1)*\lncw})
          rectangle ++(\lncw,-\lncw);
      }
    }
    \draw ({#1},{#2}) rectangle ++({10*\lncw},{-4*\lncw});
  }
  % (a) time-sharing: thread 1 for five cycles, then thread 2
  \node[anchor=south west, inner sep=1pt] at (0,0.05)
    {\iflangja{(a) 時分割}{(a) time-sharing}};
  \lnslots{0}{0}{1/3/0,2/2/0,3/3/0,4/0/0,5/0/0,6/0/3,7/0/1,8/0/2,9/0/0,10/0/2}
  \node[anchor=north, text=ln-muted] at (1.2,-0.98) {\iflangja{コア：16 命令}{core: 16 instr.}};
  % (b) CMP: one core per thread
  \node[anchor=south west, inner sep=1pt] at (3.0,0.05) {(b) CMP};
  \lnslots{3.0}{0}{1/3/0,2/2/0,3/3/0,4/0/0,5/0/0,6/1/0,7/1/0,8/0/0,9/4/0,10/2/0}
  \node[anchor=north, text=ln-muted] at (4.2,-0.98) {\iflangja{コア 1：16 命令}{core 1: 16 instr.}};
  \lnslots{3.0}{-1.44}{1/0/3,2/0/1,3/0/2,4/0/0,5/0/2,6/0/0,7/0/0,8/0/4,9/0/1,10/0/3}
  \node[anchor=north, text=ln-muted] at (4.2,-2.42) {\iflangja{コア 2：16 命令}{core 2: 16 instr.}};
  % (c) fine-grained multithreading: one thread per cycle
  \node[anchor=south west, inner sep=1pt] at (6.0,0.05)
    {\iflangja{(c) 細粒度 MT}{(c) fine-grained MT}};
  \lnslots{6.0}{0}{1/3/0,2/0/3,3/2/0,4/0/1,5/3/0,6/0/2,7/0/0,8/1/0,9/0/2,10/1/0}
  \node[anchor=north, text=ln-muted] at (7.2,-0.98) {\iflangja{コア：18 命令}{core: 18 instr.}};
  % (d) SMT: both threads in the same cycle
  \node[anchor=south west, inner sep=1pt] at (9.0,0.05) {(d) SMT};
  \lnslots{9.0}{0}{1/2/2,2/1/1,3/2/1,4/2/2,5/1/0,6/0/2,7/0/0,8/1/0,9/1/3,10/0/1}
  \node[anchor=north, text=ln-muted] at (10.2,-0.98) {\iflangja{コア：22 命令}{core: 22 instr.}};
  % axis hints
  \draw[->, ln-muted] (6.0,-1.7) -- (8.4,-1.7)
    node[midway, below, inner sep=1.5pt] {\iflangja{サイクル}{cycles}};
  \draw[->, ln-muted] (-0.12,0) -- (-0.12,-0.96)
    node[midway, rotate=90, above, inner sep=1.5pt] {\iflangja{スロット}{slots}};
  % legend
  \begin{scope}[xshift=9.0cm, yshift=-1.6cm]
    \draw[ln-rule, fill=ln-accent!60] (0,0) rectangle ++(0.24,-0.24);
    \node[anchor=west] at (0.3,-0.125) {\iflangja{スレッド 1}{thread 1}};
    \draw[ln-rule, fill=ln-tip!45] (0,-0.4) rectangle ++(0.24,-0.24);
    \node[anchor=west] at (0.3,-0.525) {\iflangja{スレッド 2}{thread 2}};
    \draw[ln-rule, fill=white] (0,-0.8) rectangle ++(0.24,-0.24);
    \node[anchor=west] at (0.3,-0.925) {\iflangja{空きスロット}{empty slot}};
  \end{scope}
\end{tikzpicture}

  \caption{Schematic use of the four slots per cycle of a core by two
    threads over ten cycles.  Each column is a cycle and each row a slot.
    Alone, each thread executes 16 instructions in these ten cycles (b);
    (a) switches threads after five cycles, (c) alternates between the
    threads that are not waiting.}
  \label{fig:l18-mt-timing}
\end{figure}

\begin{description}
  \item[Time-sharing] One core runs one thread at a time.  Together the
    threads complete no more work than a single thread, in practice
    slightly less, because every context switch costs time (not drawn in
    the figure).  The cost is that of one core.
  \item[Chip multiprocessor] Each thread runs on a core of its own.  A
    multicore processor used in this way is a
    \term{chip multiprocessor}[チップマルチプロセッサ]%
    \index{CMP|see{chip multiprocessor}} (CMP).  Both threads run at full
    speed, so the throughput doubles, but so do chip area and power, and
    each core still leaves as many slots empty as a single thread does.
  \item[Fine-grained multithreading] One core with a separate set of
    registers for each thread changes threads from cycle to cycle, so a
    cycle in which one thread must wait is used by the other.  The core is
    busier, and the only additional hardware is the small per-thread state,
    about \SI{5}{\percent} of the cost of a core.  The empty slots within a
    cycle, however, stay empty.
  \item[SMT] One core mixes instructions of both threads in the same cycle
    and fills many of the slots that a single thread would leave
    empty.\sidenote{SMT was proposed in 1995 \cite{tullsen1995}.  Intel
    shipped it as Hyper-Threading in the Pentium~4 in 2002, with two threads
    per core---still the usual x86 number---while IBM's POWER8 and POWER9
    run up to eight threads on one core.  Intel then dropped SMT again from
    the Lion Cove cores of its 2024 client chips, judging the area better
    spent on more cores.}  The
    throughput rises above that of fine-grained multithreading, though it
    stays below that of two cores, because the threads share a single set of
    resources.  SMT needs more additional hardware than fine-grained
    multithreading, about \SI{10}{\percent} of the cost of a core
    (\cref{sec:l18-smt-hw}).
\end{description}

How much SMT gains depends on the threads.  When their instructions compete
for the same execution units, the instructions of one thread find few slots
left free by the other and the gain is small; SMT approaches the throughput
of a CMP only when the resource demands of the two threads complement each
other.

\begin{example}
  \label{ex:l18-smt-cmp}
  Let SMT add \SI{10}{\percent} to the cost of a core and raise its
  throughput on two threads by \SI{35}{\percent}.  A CMP of two cores
  without SMT doubles both throughput and cost.  Relative to one core that
  time-shares the threads, the SMT core delivers $1.35/1.10\approx 1.23$
  times the throughput per unit of cost and the CMP $2/2=1$, although the
  absolute throughput of the CMP is higher.  If the two threads competed for
  the same execution units and SMT raised the throughput by only
  \SI{5}{\percent}, the SMT core would deliver $1.05/1.10\approx 0.95$, and
  two cores would be the better use of the hardware.
\end{example}

\begin{keypoint}
  A CMP gives the highest throughput, at twice the cost.  SMT gives a
  smaller gain for about \SI{10}{\percent} more cost.  When the resource
  demands of the threads complement each other, SMT is the more efficient
  use of the hardware; when the threads compete for the same resources, two
  cores can be the better choice.
\end{keypoint}

\section{Hardware changes for SMT}
\label{sec:l18-smt-hw}

SMT can be added to the out-of-order core of Lessons 7 and 8 with few
changes (\cref{fig:l18-smt-hw}).\source{Lesson 18, video 18; H\&P ch.~3;
\cite{tullsen1995}.}

\begin{figure}[htbp]
  \centering
  % An out-of-order core with SMT support and separate ROBs: per-thread PC,
% RAT, ROB and architectural registers; everything else is shared.
\begin{tikzpicture}[x=1cm,y=1cm, font=\footnotesize\sffamily, >={Stealth[length=1.8mm]},
  sh/.style={draw, fill=ln-box, minimum height=0.75cm, minimum width=1.5cm, inner sep=2pt, align=center},
  pt/.style={draw, fill=ln-accent!15, minimum height=0.36cm, minimum width=1.2cm, inner sep=1pt, font=\scriptsize\sffamily}]
  % ---- front end
  \node[pt, minimum width=0.9cm] (pc1) at (0,0.22) {PC 1};
  \node[pt, minimum width=0.9cm] (pc2) at (0,-0.22) {PC 2};
  \draw[fill=white] (0.85,0.45) -- (1.15,0.3) -- (1.15,-0.3) -- (0.85,-0.45) -- cycle;
  \draw[->] (pc1.east) -- (0.85,0.22);
  \draw[->] (pc2.east) -- (0.85,-0.22);
  \node[sh] (fetch) at (2.3,0) {\iflangja{フェッチ}{fetch}};
  \node[sh] (dec) at (4.4,0) {\iflangja{デコード}{decode}};
  \draw[->] (1.15,0) -- (fetch.west);
  \draw[->] (fetch.east) -- (dec.west);
  \node[pt] (rat1) at (6.6,0.22) {RAT 1};
  \node[pt] (rat2) at (6.6,-0.22) {RAT 2};
  \node[draw, fit=(rat1)(rat2), inner sep=2pt] (ren) {};
  \node[anchor=south, font=\scriptsize] at (ren.north) {\iflangja{リネーム}{rename}};
  \draw[->] (dec.east) -- (ren.west);
  % ---- back end
  \node[sh] (rs) at (1.1,-1.9) {RS};
  \node[sh] (ex) at (3.2,-1.9) {\iflangja{実行}{execute}};
  \node[pt, minimum width=1.6cm] (rob1) at (5.6,-1.68) {ROB 1};
  \node[pt, minimum width=1.6cm] (rob2) at (5.6,-2.12) {ROB 2};
  \node[draw, fit=(rob1)(rob2), inner sep=2pt] (robs) {};
  \node[pt] (arf1) at (8.0,-1.68) {\iflangja{レジスタ 1}{ARF 1}};
  \node[pt] (arf2) at (8.0,-2.12) {\iflangja{レジスタ 2}{ARF 2}};
  \draw[->] (ren.south) -- ++(0,-0.35) -| (rs.north);
  \draw[->] (ren.south) -- ++(0,-0.35) -| (robs.north);
  \draw[->] (rs.east) -- (ex.west);
  \draw[->] (ex.east) -- (robs.west |- ex.east) node[midway, above, font=\scriptsize] {CDB};
  \draw[->] (ex.south) -- ++(0,-0.4) -| (rs.south);
  \draw[->] (robs.east |- rob1.east) -- (arf1.west);
  \draw[->] (robs.east |- rob2.east) -- (arf2.west);
  \node[anchor=north, font=\scriptsize] at (6.95,-2.3) {\iflangja{コミット}{commit}};
  % ---- legend
  \draw[fill=ln-accent!15] (0,-3.05) rectangle ++(0.34,-0.25);
  \node[anchor=west, font=\scriptsize] at (0.4,-3.175) {\iflangja{スレッドごと}{one per thread}};
  \draw[fill=ln-box] (2.6,-3.05) rectangle ++(0.34,-0.25);
  \node[anchor=west, font=\scriptsize] at (3.0,-3.175) {\iflangja{スレッド間で共有}{shared by the threads}};
\end{tikzpicture}

  \caption{An out-of-order core with SMT for two threads and separate ROBs.
    With a shared ROB, only the program counter, the RAT and the
    architectural registers are kept per thread.}
  \label{fig:l18-smt-hw}
\end{figure}

\begin{description}
  \item[Fetch] Each thread needs its own PC.  The fetch stage takes
    instructions from both PCs and marks each fetched instruction with the
    thread it belongs to.
  \item[Decode] Decoding does not depend on the thread and is unchanged.
  \item[Rename] The renaming logic is unchanged, but each thread needs its
    own RAT: register R1 of thread 1 and register R1 of thread 2 are
    different architectural registers.
  \item[Dispatch and execution] Unchanged.  The reservation stations, the
    execution units and the CDB are shared; after renaming, operands are
    identified by tags that are unique across both threads, so instructions
    of different threads cannot be confused.
  \item[ROB] Each thread can have a separate ROB with its own commit point.
    The ROB is a large structure, however, so duplicating it is expensive.
    More often in practice, the instructions of both threads are interleaved
    in a single shared ROB, each entry recording the thread of its
    instruction.\margin{The Pentium~4 partitioned rather than shared: while
    both threads ran, each got half of the 126 ROB entries
    \cite{marr2002}.}  The small price is that a completed instruction of one
    thread may have to wait for older instructions of the other thread to
    commit.
  \item[Commit] Each thread needs its own architectural register file (ARF),
    into which its instructions commit.
\end{description}

With a shared ROB, the duplicated structures---PC, RAT and ARF---are small;
the expensive parts of the core---ROB, reservation stations, execution
units, caches and branch predictors---are shared.  This is why SMT adds
only about \SI{10}{\percent} to the cost of a core.\margin{Intel put
Hyper-Threading on the Pentium~4 at about \SI{5}{\percent} more die area for
\SI{15}{\percent}--\SI{30}{\percent} more throughput \cite{marr2002}.}

\section{SMT, the data cache and the TLB}
\label{sec:l18-smt-tlb}

The threads that share an SMT core may belong to different processes, each
with its own virtual address space.\source{Lesson 18, video 19; H\&P
Appendix B.}  The same virtual address then refers to different data in each
thread, and the caches and the TLB of Lessons 13 and 14 must not confuse the
two.
\begin{itemize}
  \item A \term*{virtually accessed cache} is indexed and tagged with the
    virtual address alone.  A block cached for one thread produces a hit for
    the other thread's access to the same virtual address and returns the
    wrong data.  Such a cache can be used with SMT only if each entry also
    records its thread and a hit requires the thread to match as well.
  \item A \term*{VIPT cache} and a physically accessed
    cache compare physical tags.  Two threads hit the same block only when
    they access the same physical address, and then they really do access
    the same data.  These caches work correctly with SMT as long as the TLB
    delivers the right translation for each thread.
  \item The \term*{TLB} must therefore be thread-aware: each entry records the
    thread it belongs to, the thread identifier of the access is passed into
    the lookup, and a hit requires both the thread and the virtual page
    number to match (\cref{fig:l18-smt-tlb}).\caution{What must differ is the
    address space, not the thread: two threads of one process share their
    translations.  Real TLBs therefore tag entries with an address-space
    identifier---PCID on x86-64, ASID on Arm.}
\end{itemize}

\begin{figure}[htbp]
  \centering
  % Thread-aware TLB of an SMT core: every entry records the thread it
% belongs to, and a lookup must match both the thread and the VPN.
\begin{tikzpicture}[x=1cm,y=1cm, font=\footnotesize\sffamily, >={Stealth[length=1.8mm]},
  cell/.style={draw, minimum height=0.42cm, inner sep=1pt, font=\footnotesize\ttfamily},
  hd/.style={inner sep=1pt, font=\scriptsize\sffamily, text=ln-muted}]
  % ---- lookup key
  \node[cell, fill=ln-accent!15, minimum width=0.8cm] (kt) at (0,0.21) {1};
  \node[cell, fill=ln-accent!15, minimum width=1.2cm, anchor=west] (kv) at (kt.east) {0x12};
  \node[hd, anchor=south] at (kt.north) {\iflangja{スレッド}{thread}};
  \node[hd, anchor=south] at (kv.north) {VPN};
  % ---- TLB entries: thread | VPN | frame
  \foreach \r/\t/\v/\f in {0/0/0x12/0x3A, 1/1/0x12/0x07, 2/0/0x40/0x15, 3/1/0x2C/0x51} {
    \ifnum\r=1 \def\lnfill{ln-tip!25}\else \def\lnfill{white}\fi
    \node[cell, fill=\lnfill, minimum width=0.8cm] (t\r) at (4.2,{0.63-\r*0.42}) {\t};
    \node[cell, fill=\lnfill, minimum width=1.2cm, anchor=west] (v\r) at (t\r.east) {\v};
    \node[cell, fill=\lnfill, minimum width=1.2cm, anchor=west] (f\r) at (v\r.east) {\f};
  }
  \node[hd, anchor=south] at (t0.north) {\iflangja{スレッド}{thread}};
  \node[hd, anchor=south] at (v0.north) {VPN};
  \node[hd, anchor=south] at (f0.north) {\iflangja{フレーム}{frame}};
  \node[font=\footnotesize, anchor=north] at (v3.south) {TLB};
  % ---- lookup and result
  \draw[->] (kv.east) -- (t1.west |- kv.east)
    node[midway, above, font=\scriptsize, align=center]
    {\iflangja{スレッドと VPN が}{thread and VPN}\\\iflangja{ともに一致}{must both match}};
  \node[cell, fill=ln-tip!25, minimum width=1.2cm] (res) at (9.2,0.21) {0x07};
  \node[hd, anchor=south] at (res.north) {\iflangja{物理フレーム}{physical frame}};
  \draw[->] (f1.east) -- (res.west);
  \node[hd, anchor=north west, align=left] at (3.8,-1.35)
    {\iflangja{同じ VPN 0x12 でもスレッド 0 はフレーム 0x3A に対応する}{VPN 0x12 of thread 0 maps to frame 0x3A instead}};
\end{tikzpicture}

  \caption{A thread-aware TLB.  The same virtual page number is translated
    to different frames for the two threads.}
  \label{fig:l18-smt-tlb}
\end{figure}

\section{SMT and cache performance}
\label{sec:l18-smt-cache}

The caches of a core are shared by all its SMT threads, which has one
advantage and one serious drawback.\source{Lesson 18, videos 20--22; H\&P
ch.~3.}
\begin{itemize}
  \item \textbf{Fast data sharing.}  When one thread stores a value that
    the other thread then loads, the load finds the data already in the
    cache.
  \item \textbf{Shared capacity.}  The threads also share the capacity and
    the associativity of the cache, and so they compete for it.
\end{itemize}

The competition is described in terms of the
\term{working set}[ワーキングセット] of a thread: the set of blocks it is actively using
at a given time.  Let $\mathit{WS}_0$ and $\mathit{WS}_1$ be the sizes of
the working sets of two SMT threads and $\mathit{WS}_{01}$ the size of the
part both use.  As long as
\begin{equation}
  \mathit{WS}_0 + \mathit{WS}_1 - \mathit{WS}_{01} \;\le\; \text{cache size},
  \label{eq:l18-ws}
\end{equation}
both threads find their blocks in the cache almost as often as when running
alone.  When the combined working set exceeds the cache, the threads keep
evicting each other's blocks and the cache \emph{thrashes}.  If each working
set would fit on its own, SMT can then be slower than running the threads
one at a time: a time-shared thread has the whole cache to itself for
millions of cycles and refills it only after each context switch, whereas
SMT threads evict each other's blocks all the time.\caution{SMT raises
throughput only if the cache holds the working sets of all threads at the
same time.}  A scheduler should therefore run together threads whose
combined working set fits the cache, and run a thread whose working set
alone exceeds the cache without SMT partners: it would evict their blocks
without gaining anything itself.

\begin{example}
  A core has a \SI{64}{KB} L1 data cache.\margin{Real L1 data caches
  are smaller: 32~KB on AMD's Zen~4, 48~KB on Intel's Golden Cove, with
  1--2~MB of L2 per core behind them.}  Threads A, B and C have working
  sets of \SI{40}{KB}, \SI{36}{KB} and \SI{28}{KB}; each fits into the cache
  alone.  A and B work on common data, \SI{20}{KB} of which appears in both
  working sets, so running them together with SMT needs
  $40+36-20=\SI{56}{KB}$, which fits.  A and C share no data and together
  need \SI{68}{KB}; with SMT they thrash, and time-sharing them may well be
  faster.
\end{example}

\begin{keypoint}[Lesson summary]
  Multiprocessors are MIMD machines.  They have replaced ever-faster
  uniprocessors because wider pipelines give diminishing returns and higher
  frequencies cost power roughly as $f^{3}$, whereas replicated cores turn
  Moore's-law transistors into throughput---but only for parallel programs,
  whose speedup levels off (\cref{eq:l18-amdahl-cores}).  Centralized shared
  memory (UMA, SMP) is simple but limited to a few cores by memory size and
  bandwidth; distributed shared memory (NUMA) scales to large machines at the price
  of slow remote accesses.  Message passing is hard to program correctly,
  shared memory hard to make fast.  A core can run several threads by
  time-sharing, by coarse- or fine-grained multithreading, or by SMT, which
  duplicates little more than the PC, the RAT and the ARF, needs a
  thread-aware TLB, and pays off only if the shared cache holds the combined
  working set (\cref{eq:l18-ws}).  Cores that share memory must also keep
  their caches consistent, the topic of Lesson 19.
\end{keypoint}

\end{document}
